\chapter{Code availability}\label{app:code}

\begin{table}[H]
	\centering
	\captionsetup{width=1.0\textwidth} 
	\begin{tabularx}{\textwidth}{>{\raggedright\arraybackslash}p{5.2cm} >{\raggedright\arraybackslash}X}
		\toprule
		\textbf{Repository} & \textbf{Function}\\
		\midrule
		\textbf{\href{https://github.com/yourusername/VisualFNS.jl}{\texttt{VisualFNS.jl}}}
		& Producing simulations, visualisations, and raw data of gaze trajectories through numerically integrating the system \cref{eq:fns_system} with kernel adaptation under \cref{eq:kernel_law}. Important techniques in solving this system were the use of the Euler–Maruyama numerical scheme for SDE integration and the Fourier-space spectral solution of kernel updates.\\
		\addlinespace
		\textbf{\href{https://github.com/yourusername/VisualFNS_Analysis.jl}{\texttt{VisualFNS\_Analysis.jl}}}
		& Analysing trajectory data which was output from \href{https://github.com/yourusername/VisualFNS.jl}{\texttt{VisualFNS.jl}}. Involved producing power spectra, marking intervals between saccade-like transitions, and analysing superdiffusive scaling.\\
		\addlinespace
		\textbf{\href{https://github.com/yourusername/VisualData_Analysis.jl}{\texttt{VisualData\_Analysis.jl}}}
		& Loading in the data provided from \textcite{naher_primate_saccades}, and analysing it in the same way as in \href{https://github.com/yourusername/VisualFNS_Analysis.jl}{\texttt{VisualFNS\_Analysis.jl}}. Involved producing power spectra, marking intervals between saccade-like transitions, and analysing superdiffusive scaling. Originally forked from \href{https://github.com/yourusername/FractionalSaccades.jl}{\texttt{FractionalSaccades.jl}}.\\
		\addlinespace
		\textbf{\href{https://github.com/yourusername/FractionalNeuralSampling.jl}{\texttt{FractionalNeuralSampling.jl}}}
		& The project which provides basic methods to formalise the FNS equation and its variants in the syntax solvable through \texttt{DifferentialEquations.jl}.\\
		\bottomrule
	\end{tabularx}
	\caption{\textbf{A list of repositories containing the code used in this thesis.} All code was written in Julia, apart from some helper scripts written in Python.}
	\label{tab:code_availability}
\end{table}